\section{Weak-alpha weak-\(n\) and C1: the complete parent}
\label{sec:weak-alpha-weak-n-c1}

\paragraph{Scope and retained inputs.}
The source is an actual physical canonical finite noncommuting contact,
all coefficients are bounded in the source field, and \(d=1\).
Either it is C2 with \(\lambda=\beta=0\) and \(n,\alpha\) infinitesimal,
or it is C1 with \(\lambda=n=0\), \(\alpha\) infinitesimal and arbitrary
bounded \(\beta\geq0\). All strict source gaps retain their original
meaning. We prove the tensor-SOS statement for the entire parent at
every independent projective testing angle and every state.
No prior-complement exclusion is assumed.
The common source-entry proof WA1--WA10 is used on precisely its
weak-\(n\) domain. In particular \(C,A,D,r,y,r-y,1+y-r\) are positive
units, \(x\) is bounded, and \(x_0=-r_0\).
Boundedness, infinitesimality and residue refer
to the separated bounded tensor ring and its real residue map.
The strict-Gram lemma is applied only after a finite explicit row
allocation; no mixed tensor is inverted.

Use the full quarter-parent WA12 and its exact source coordinates
\[
 p=v_c,\quad q=Z_c>0,\quad P=p^2,\quad Q=q^2,\qquad
 W_c=1,\quad u_c=rpq,\quad X_c=xpq^2,\quad Y_c=ypq.
 \tag{WN1}
\]
Thus \(P,PQ\) are positive infinitesimals. Put
\[
 H_0=m+\beta+\tfrac12,\quad H_X=H_0+n,\quad
 K=C+D+\tfrac12,\quad H_Y=K+A.
 \tag{WN2}
\]
These are positive units. The exact scalar coefficients
\(\mathcal C,E,\mathcal F,G,J,\mathcal K\), the contact balances, and
the complete first normal polynomials are those displayed in
WA14, WA16 and WA19; those identities retain \(\beta\) and \(n\).
Their relevant residues are
\[
 \mathcal C_0=\frac{r_0(1+y_0+4m_0y_0)}2>0,\qquad
 \mathcal F_0=\frac{r_0^2(1+y_0)}2>0,
 \qquad G,J,\mathcal K\geq0 .
 \tag{WN3}
\]
Both \(\mathcal C,\mathcal F\) are consequently positive units.
The exact source \(X\)-jet gives, with the WA19 coefficients,
\[
\begin{gathered}
 c_0=H_0(r+x)+nx,\quad c_2=nr^2x,\quad
 c_3=(D+n+\tfrac12)r^2x,\\
 c_0=m/Q-Pc_1-PQc_2-P^2Qc_3,\\
 E=\alpha r^2+c_0^2/H_X+nH_0r^2/H_X .
\end{gathered}\tag{WN4}
\]
The last equality is an exact completion inside the source field.

\subsection{Two finite positive circuit identities}

For \(V=v^2,T=z^2\) put
\[
 R=z(V-P),\quad S=v(T-Q),\quad
 \Psi=V^2T^2+PQ^2V+P^2QT-3PQVT .
\]
Then \(\Psi=QR^2+PS^2+vzRS\). It is an explicit six-square
cubic circuit: if \(L_1^3=V^2T^2,L_2^3=PQ^2V,L_3^3=P^2QT\),
then \(L_1L_2L_3=PQVT\), and
\[
 \Psi=\sum_{i=1}^3 L_i
 \bigl\{(L_1-L_2)^2-(L_1-L_2)(L_1-L_3)+(L_1-L_3)^2\bigr\}.
 \tag{WN5}
\]
The real \(2\)-by-\(2\) Gram in each brace is strictly positive.
All radicals in this formula are elementary separated monomials:
source radicals and testing radicals are taken in their own real
closed fields and then tensored.

The balanced degree-\((3,3)\) term has the polynomial identity
\[
 a^{10}+2b^{10}+2c^{10}-5a^2b^4c^4
 =2(b^5-c^5)^2
 +(a^2-bc)^2(a^6+2a^4bc+3a^2b^2c^2+4b^3c^3).
 \tag{WN6}
\]
For nonnegative elementary monomials \(a,b,c\), the bracket is the
sum of squares of the four rows
\[
 a^3,\qquad \sqrt2a^2\sqrt{bc},\qquad
 \sqrt3abc,\qquad 2(bc)^{3/2}.
\]
This is a five-square separated certificate, including vanishing
variables. Applied with
\[
 a^{10}=V^3T^3,\quad b^{10}=P^2Q^3V,\quad c^{10}=P^3Q^2T,
\]
it proves positivity of the entire \(B_{33}\) in WA17, rather than
discarding a signed high-order remainder.

A further exact identity supplies any fixed positive real pivot
on the row \(PS\):
\[
 2P\Psi+VT R^2-(PS)^2=(PS+vzR)^2+2PQR^2.
 \tag{WN7}
\]
Indeed, for every fixed \(b>0\),
\[
\begin{split}
 \mathcal C R^2+\mathcal F\Psi={}&
 (\mathcal C-bVT+2bPQ)R^2+(\mathcal F-2bP)\Psi\\
 &+b(PS)^2+b(PS+vzR)^2.
\end{split}\tag{WN8}
\]
On the charts where this is used, \(VT,P,PQ\) are infinitesimal,
so the original two unit residues are unchanged.

\subsection{Bounded source ratio \(Q\)}

From WA5, WN4 and bounded \(Q\),
\[
 m\asymp Q^2(n+P),\quad \alpha\asymp P^2,\quad
 r+x=O(n+P),\quad E\asymp n+P^2.
 \tag{WN9}
\]
In particular \(m,n,E,\alpha\) are infinitesimal. None of these
estimates contains an independent testing coordinate.

First take \(W=1,Z=z\) bounded and \(u,v\) infinitesimal. Use the
entire graph \(u=rvz,X=xvz^2,Y=yvz\). The exact scalar is
\[
 \Phi=\mathcal CR^2+ES^2+\mathcal F\Psi
          +GB_{32}+JB_{23}+\mathcal K B_{33},
 \tag{WN10}
\]
where the last three terms retain their positive full certificates.
The strict soft rows are \(R\), the six cubic rows in WN5,
\(\sqrt E\,S\), and \(PS\); WN8 supplies any needed real pivot
on the last row.

The exact monic normal divisions WA20 apply unchanged.
Their \(R\)-coefficients contain \(v\) or \(vz\), hence are
infinitesimal. Their \(S\)-coefficients are \(O(E+P)\):
in particular
\[
 a_{01}=\alpha r+H_0(r+x)=O(E+P),\qquad
 a_{12}=nr x^2,\qquad m/Q=O(E+P).
 \tag{WN11}
\]
Every \(ES\) contribution has infinitesimal coefficient \(\sqrt E\)
on \(\sqrt E S\). Coefficients on \(PS\) may be bounded units.
There are only finitely many of them, so choose \(b\) in WN8
larger than their finite real Schur cost. This uses a real constant,
not an inverse of a mixed tensor.

The three hard normals
\(N_u=u-rvz,N_X=X-xvz^2,N_Y=Y-yvz\) are infinitesimal.
Their complete quadratic residue is
\[
 K_0(N_u-N_Y)^2+A_0N_Y^2+
          (\beta_0+\tfrac12)(N_X+\bar zN_u)^2\succ0.
 \tag{WN12}
\]
All higher-normal terms have degree at least three and bounded
coefficients, hence allocate as infinitesimal Gram entries on these
hard rows. WN10--WN12 prove the full parent in this chart.

For the complementary local states take \(Z=1,W=w\) infinitesimal,
set \(a=u/r\), and use
\[
 v=aw,\quad X=xa,\quad Y=yaw,\qquad U=a^2,\quad T=w^2.
 \tag{WN13}
\]
The entire scalar is
\[
\begin{split}
 \Phi'={}&\alpha+mUT^2+EU+hUT+\mathcal C U^2T^2
       +\mathcal F U^2T+GU^3T^2+JU^2+\mathcal K U^3T .
\end{split}\tag{WN14}
\]
Here \(a=(1/r)\otimes u\) is elementary; its roots remain separated.
The strict rows are \(P,a^2w,\sqrt E\,a\).
The term \(hUT\) is harmless: the contact balances give
\(h=O(P+E)\), and writing \(h=Ph_P+Eh_E\) with bounded source
coefficients gives the infinitesimal Gram terms
\(h_Pw\,P(a^2w)+h_Ew^2(\sqrt E a)^2\).
All other terms of WN14 are positive.

For completeness, half the entire \(v,Y,X\) derivatives have
the polynomial forms
\[
\begin{aligned}
 L_v'&=aw(d_0+mT+d_1U+d_2UT+d_3U^2+d_4U^2T),\\
 L_Y'&=aw(e_0+e_1U+e_2UT+e_3U^2T),\\
 L_X'&=a(c_0-mT+c_2U+c_1UT+c_3U^2T).
\end{aligned}\tag{WN15}
\]
Every coefficient is bounded. Source jets at
\(U=PQ,T=1/Q\) give
\[
\begin{split}
d_0&=-m/Q-d_1PQ-d_2P-d_3P^2Q^2-d_4P^2Q=O(n+P),\\
e_0&=-e_1PQ-e_2P-e_3P^2Q=O(P).
\end{split}
\]
Also \(c_0,m=O(n+P)\) and \(c_2=O(n)\).
Splitting \(n+P=O(E+P)\) in these displayed complete polynomials
allocates all first normals with infinitesimal coefficients to
\(P,a^2w,\sqrt E a\). For example \(awP\), \(awE\), and
\(a^3w\) have coefficients \(aw\), \(w\sqrt E\), and \(a\)
on those three rows; the other monomials have still more small
testing factors. The complete hard residue is
\[
 (\beta_0+\tfrac12)N_X^2+
 D_0N_v^2+2D_0N_vN_Y+H_{Y,0}N_Y^2\succ0,
 \tag{WN16}
\]
where \(N_v=v-aw,N_Y=Y-yaw,N_X=X-xa\).
The higher-normal argument is again exact polynomial allocation.
This proves all local states when \(Q\) is bounded.

\subsection{Unbounded source ratio: the \(W\)-chart}

Suppose \(Q\) is unbounded. The source estimates and positive
contact balances give
\[
 PQ=o(1),\quad m\asymp Q^2(n+P),\quad nQ^2,PQ^2=O(1),
 \quad \alpha\asymp P^2Q .
 \tag{WN17}
\]
WN4 now gives
\[
 c_0\sim m/Q,\quad r+x\sim m/(H_XQ),\quad
 E\asymp n+(m/Q)^2,\quad h/m=o(1).
 \tag{WN18}
\]
To justify the first comparison before its use, divide the three
errors \(Pc_1,PQc_2,P^2Qc_3\) by \(m/Q\).
Their absolute values are \(O(1/Q),O(P),O(P)\), respectively;
\(n/(m/Q)=O(1/Q)\). The positive term
\(\alpha/(m/Q)^2=O(1/Q)\) follows from
\(m\geq cPQ^2\). The last comparison follows term by term from
WA16 using these same inequalities.

Take \(W=1,Z=z\) bounded, \(u,v\) infinitesimal. If \(m\) is a
unit, retain the entire \(n,\alpha\) packets. The remainder lies
in the square of \((u,v,X,Y)\). Its real quadratic form is
\[
\begin{split}
 &H_0\{X+zu-(m/H_0)v\}^2+
       \frac{m(\beta+1/2)}{H_0}v^2\\
 &\qquad+H_Y\{Y-(Ku-Dzv)/H_Y\}^2
       +a_*(u-rzv)^2,\qquad a_*=AK/H_Y>0,
\end{split}\tag{WN19}
\]
with coefficients interpreted in residue. WA10 verifies this
identity. It is strict for every bounded \(z\), so the complete
higher terms give a strict Gram.

If \(m\) is weak, use the graph WA13. Its entire scalar has strict
rows \(\sqrt m\,v,v^2z,\sqrt\alpha\,z\), since
\(mV+hVT=mV[1+(h/m)T]\), while \(\mathcal C,\mathcal F,E,G,J,\mathcal K\)
are nonnegative. The complete source jet equations yield
\[
 a_{00}=O(m),\quad e_{00}=O(PQ),\quad
 a_{01},c_0=O(m/Q).
 \tag{WN20}
\]
Thus \(a_{00}/\sqrt m\), \(e_{00}/\sqrt m\),
\(a_{01}/\sqrt m\), \(c_0/\sqrt m\) are infinitesimal.
For \(e_{00}\), use \(PQ/\sqrt m=O(\sqrt P)\).
Every other monomial of the complete normals WA19 has a factor
\(v^3z\) or \(v^3z^2\), so it has an infinitesimal coefficient
on \(v^2z\). The hard residue is WN12; all higher terms are
allocated exactly as there. This proves the whole \(W\)-chart.

\subsection{Unbounded source ratio: the complementary scalar}

Use WN13 with \(a,w\) infinitesimal. Introduce only source quantities
\[
\begin{gathered}
 \epsilon=Q^{-1}=o(1),\quad \Lambda=PQ=o(1),\quad
 \tau=\Lambda/\epsilon=PQ^2=O(1),\\
 \nu=n/\epsilon^2=O(1),\qquad N=E/\epsilon^2=O(1).
\end{gathered}\tag{WN21}
\]
The quantities \(\epsilon,\Lambda,\tau,N\) are positive; \(\nu\geq0\).
Let
\[
 A_1=a(T-\epsilon),\quad K_1=UT-\Lambda\epsilon,\quad
 \mathcal M=U^2T+\tau UT^2+\Lambda^2\epsilon-3\Lambda UT .
 \tag{WN22}
\]
The whole scalar WN14, using exactly the three contact balances, is
\[
\begin{split}
 \Phi'={}&\mathcal C K_1^2+NA_1^2+\mathcal F\mathcal M\\
 &+G\{2\Lambda K_1^2+UT^2(U-\Lambda)^2\}\\
 &+J\{(U-\Lambda)^2+2\Lambda A_1^2/\epsilon^2\}
       +\mathcal K\mathcal B_{33},\\
 \mathcal B_{33}={}&U^3T+2(\Lambda^2/\epsilon)UT^2
                      +2\Lambda^3\epsilon-5\Lambda^2UT .
\end{split}\tag{WN23}
\]
The apparent \(1/\epsilon^2\) in the positive \(J\) packet is
a source-field quantity, and \(J/\epsilon^2=\nu r^2x^2\) is bounded.
All positive terms are retained. WN6 applies to \(\mathcal B_{33}\)
with \(a^{10}=U^3T,b^{10}=(\Lambda^2/\epsilon)UT^2,
c^{10}=\Lambda^3\epsilon\).
Here \(\Lambda^2/\epsilon=\tau\Lambda\) is bounded.

Write \(a=t^3,w=s^3,\epsilon=h^6,\tau=c^6\), taking positive
elementary roots. Then \(t,s,h\) are infinitesimal and \(c\)
is bounded; it is allowed to be a unit. Put
\[
\begin{gathered}
 L_1=t^4s^2,\quad L_2=c^2t^2s^4,\quad L_3=h^6c^4,\\
 f_1=t^2s,\quad f_2=cts^2,\quad f_3=h^3c^2,\\
 R_{2i-1}=f_i(L_1-L_2),\qquad R_{2i}=f_i(L_1-L_3)
 \quad (i=1,2,3).
\end{gathered}\tag{WN24}
\]
Then \(\mathcal M=\sum_{i=1}^3(R_{2i-1}^2-R_{2i-1}R_{2i}+R_{2i}^2)\),
a strict six-row constant Gram. The following are exact identities:
\[
\begin{aligned}
 aw(U-\Lambda)&=t^3R_1+c^2ts^2R_2,\\
 a^2w(T-\epsilon)&=-h^6R_1+s^6R_2-h^3s^3R_5,\\
 \tau aw(T-\epsilon)&=c^2ts^2(R_2-R_1)-c^3s^3R_3,\\
 \tau\sqrt\epsilon\,A_1&=t^3R_6-t(t^2+c^2s^2)R_5 .
\end{aligned}\tag{WN25}
\]
Every coefficient on the right is infinitesimal, even when \(c\)
is a unit. These are finite linear conductors, not limiting
approximations. The further quadratic conductor is
\[
 \boxed{\ \Lambda(U-\Lambda)K_1
       =\Lambda\mathcal M-(\tau\sqrt\epsilon\,A_1)^2.\ }
 \tag{WN26}
\]
Its right side is an infinitesimal Gram on the six cubic rows.
It remains so after multiplication by any bounded signed coefficient.

If \(\tau\) is a unit, then \(K_1\) itself has an infinitesimal
linear representation on those rows:
\[
 K_1=\tau^{-1/3}
       \{\sqrt{L_1}(R_2-R_1)+\sqrt{L_3}R_6\}.
 \tag{WN27}
\]
This follows by expansion as
\(\tau^{-1/3}(L_1L_2-L_3^2)=UT-\Lambda\epsilon\).
Consequently any fixed real pivot on \(K_1\) can be supplied at
infinitesimal cost to the strict \(\mathcal M\) Gram in this case.

\subsection{Every complementary first normal}

Let \(D_U,E_U\) denote the monic \(U\)-divided differences of
the bracket polynomials in the first two lines of WN15.
The source jets at \(U=\Lambda,T=\epsilon\) give the exact
identities
\[
\begin{aligned}
 L_v'={}&aw\{D_U(U,T)(U-\Lambda)
                 +(m+\Lambda d_2+\Lambda^2d_4)(T-\epsilon)\},\\
 L_Y'={}&aw\{E_U(U,T)(U-\Lambda)
                 +(\Lambda e_2+\Lambda^2e_3)(T-\epsilon)\},\\
 L_X'={}&-mA_1+
       a(c_2+c_3\Lambda\epsilon)(U-\Lambda)
                  +a(c_1+c_3U)K_1 .
\end{aligned}\tag{WN28}
\]
There are no omitted monomials. Define
\[
 d_*=(c_2+c_3\Lambda\epsilon)/\epsilon
       =\nu\epsilon r^2x+\Lambda c_3 .
\]
Since \(\epsilon(U-\Lambda)=K_1-U(T-\epsilon)\), the last line is
\[
 L_X'=-mA_1+P_X,\qquad
 P_X=a(c_1+c_3U+d_*)K_1-d_*UA_1 .
 \tag{WN29}
\]
The exact source equations also give
\[
\begin{aligned}
 m&=N+\tau(\mathcal F+G\Lambda\epsilon+2\mathcal K\Lambda)
                            +2\nu r^2x^2\Lambda,\\
 c_0/\epsilon&=m-\Lambda B,\qquad
 B=c_1+\Lambda c_3+\nu\epsilon r^2x,\\
 N&=\alpha r^2/\epsilon^2+
             (m-\Lambda B)^2/H_X+\nu H_0r^2/H_X .
\end{aligned}\tag{WN30}
\]
In particular
\[
 N,\nu,\tau=O(m),\quad \Lambda=O(m\epsilon),\quad
 N=m^2/H_X+\nu H_0r^2/H_X+O(m\Lambda).
 \tag{WN31}
\]
For the error term use \(\alpha/\epsilon^2=O(\tau\Lambda)\),
which is the exact positive alpha balance WN23/WA16.
All these comparisons are inside the source field.
We split by the bounded source ratio \(\nu/m\); \(m>0\).

\paragraph{Substantial \(\nu/m\).}
If \(\nu/m\) has positive residue, keep \(S_0=\sqrt m A_1\).
The \(N_X,S_0\) Gram, including the entire \(-mA_1\) first
normal, has residue from
\[
 \begin{pmatrix}H_X&-\sqrt m\\-\sqrt m&N/m\end{pmatrix}.
 \tag{WN32}
\]
Its determinant is
\((\nu/m)H_0r^2+O(\Lambda)\), a positive unit.
The \(v,Y\) first normals in WN28 use the first conductor WN25;
their remaining coefficients on \(S_0\) are
\(O(\sqrt m\,w)\), hence infinitesimal.
In WN29 the \(K_1\)-coefficient contains \(a\), and
\(d_*U/\sqrt m=O(\sqrt m\,\epsilon U)\) is infinitesimal.
Thus all other first normals have infinitesimal coefficients
on \(K_1\), the six cubic rows, and \(S_0\).
The \(v,Y\) hard block in WN16 is strict. All higher terms have
at least three infinitesimal hard normals and bounded coefficients.
This proves the full parent on this source subdomain.

\paragraph{Thin \(\nu/m\), including literal \(n=0\).}
Assume \(\nu/m=o(1)\). Equations WN30--WN31 imply
\(\tau/m\) is a positive unit: after division by \(m\),
\[
 1=N/m+(\tau/m)\mathcal F+o(1),\qquad
 N/m=m/H_X+o(1),
\]
and \(1-m/H_X=(n+\beta+1/2)/H_X\) has positive residue.
Write
\[
 N-m^2/H_X=\nu\xi+\eta,\qquad
 \xi=H_0r^2/H_X>0,\qquad \eta=O(m\Lambda).
\]
Complete only the source coefficient \(H_X\):
\[
\begin{split}
 H_XN_X^2-2mN_XA_1+NA_1^2
 &=H_X\widetilde N_X^2+\nu\xi A_1^2+\eta A_1^2,\\
 \widetilde N_X&=N_X-(m/H_X)A_1.
\end{split}\tag{WN33}
\]
Retain the strict row \(\sqrt\nu A_1\) when \(\nu>0\);
if \(\nu=0\), all terms allocated to it vanish literally.
The signed error is an infinitesimal cubic Gram, since
\[
 m\Lambda A_1^2=(m/\tau)(\tau\sqrt\epsilon\,A_1)^2.
 \tag{WN34}
\]
The first two normals in WN28 use WN25 and bounded \(m/\tau\).
For \(P_X\), its \(K_1\)-coefficient contains \(a\), while its
two other contributions have infinitesimal allocations
\[
 \nu\epsilon U A_1=(\sqrt\nu\,\epsilon U)(\sqrt\nu A_1),
 \qquad \Lambda U A_1=(\sqrt\epsilon\,U)
                         (\tau\sqrt\epsilon\,A_1).
 \tag{WN35}
\]
Therefore \(2P_X\widetilde N_X\) has infinitesimal Gram entries.

The shift in WN33 also creates the entire scalar term
\(2(m/H_X)A_1P_X\), which must not be dropped.
Its main factor is controlled by the exact identity
\[
 mU(T-\epsilon)K_1
   =mK_1^2-(m/\tau)\Lambda(U-\Lambda)K_1 .
 \tag{WN36}
\]
The second term is infinitesimal by WN26. If \(m\) is
infinitesimal, the first is an infinitesimal \(K_1\)-Gram.
If \(m\) is a unit, then \(\tau\) is a unit and WN27
makes it an infinitesimal cubic Gram. In that case a positive
real \(K_1\)-pivot may be reinserted by WN27 at infinitesimal
cubic cost. This allocation remains valid with each bounded
factor \(c_1+c_3U+d_*\).
The remaining scalar factors are multiples of
\(m\nu\epsilon U A_1^2\) and \(m\Lambda U A_1^2\);
they have infinitesimal Gram entries on \(\sqrt\nu A_1\)
and the fourth conductor in WN25, respectively.
This accounts for the whole scalar effect of the completion.

\subsection{The entire higher-normal restoration in the thin case}

The full parent is quadratic in \(X\). Its coefficient beyond
the source constant \(H_X\) is exactly the nonnegative polynomial
\[
 D_X=(n+\beta+\tfrac12)v^2+nu^2+
                         (n+D+\tfrac12)u^2v^2 .
 \tag{WN37}
\]
Retain the entire packet \(D_XN_X^2\), evaluated at the actual
independent \(u,v\), as separated squares. Each monomial coefficient
is a nonnegative source scalar times a testing square. It remains
a positive whole packet after WN33.
The remaining polynomial has no quadratic \(N_X\) term except
\(H_XN_X^2\). At the complementary graph its half mixed
derivatives have the exact forms
\[
 [N_XN_v]=w[-m+a^2B_v],\qquad
 [N_XN_Y]=a^2wB_Y,
 \tag{WN38}
\]
where \(B_v,B_Y\) are bounded polynomials fixed by the full
parent expansion. After the shift, the new first-normal terms
are bounded multiples of \(\tau aw(T-\epsilon)\), using
\(m/\tau=O(1)\); WN25 gives infinitesimal coefficients.
Every other remaining monomial linear in \(N_X\) has degree at
least two in the other hard normals. Replacing \(N_X\) by
\(\widetilde N_X+(m/H_X)A_1\) gives an infinitesimal coefficient
on a product of at least two hard normals, since \(A_1=o(1)\).
Terms involving \(\widetilde N_X\) have total hard-normal
degree at least three and are handled in the same way.
Terms independent of \(N_X\) are unchanged. The full real hard block is
\[
 H_{X,0}\widetilde N_X^2+
       D_0N_v^2+2D_0N_vN_Y+H_{Y,0}N_Y^2\succ0.
 \tag{WN39}
\]
Equations WN23--WN39 therefore give a finite strict Gram for
the whole parent, including all first normals, every higher term,
all \(\beta\)-openings and literal \(n=0\).
This completes the double-infinitesimal testing corner for both
source ratios and all local states.

\subsection{Every other angle, all states and literal endpoints}

Retain the entire \(m,n,\alpha,\beta\) packets and call the
remainder \(F_0\). Only its residue is the critical \(C,A,D\)
source in WA10; no global physicality of the subtracted finite
source is asserted. The following are complete polynomial row
allocations, so their strict residue Grams apply to the original
nearby coefficients.

For bounded \(u,v\) with \(v\) a unit use \((uW,X,Y,Z)\).
The \(Z=0\) principal block is positive definite and
\[
 \det\operatorname{Mat}_{(uW,X,Y,Z)}\bar F_0
 =D_u^2D_v^2(q_{12}v^4+q_{21}u^2v^2)/16>0.
 \tag{WN40}
\]
Both \(q_{12},q_{21}\) are positive units. If \(u\) is a unit
and \(v\) infinitesimal, use instead
\((uW-Y,Y,X+uZ,vZ)\). The entire parent is polynomial in
these rows with bounded coefficients; the only algebraic
coefficient divisions are by the pure testing unit \(u\).
Its determinant at \(v=0\) is
\(D_u^2q_{21}u^2/16>0\), and the first three-row principal block
is strict. This covers all bounded angle pairs outside the
double-infinitesimal corner.

For Bob reciprocal \(b=1/v=o(1)\), Alice bounded, the rows
\((uW,X,Y,Z)\) give the real form
\[
 C_0(uW-Y)^2+(A_0+\tfrac12)D_uY^2+
                 D_0(uX-Z)^2+\tfrac12D_uX^2\succ0.
 \tag{WN41}
\]
For Alice reciprocal \(s=1/u=o(1)\) and Bob a bounded unit,
the original state rows give
\[
 C_0D_vW^2+A_0D_vY^2+
       (D_0+\tfrac12)(W+vX)^2+\tfrac12(vY-Z)^2\succ0.
\]
If \(s,v\) are both infinitesimal, use
\((W,Y,vX,\Gamma)\), where \(\Gamma=sX+Z\).
The entire reciprocal parent is polynomial in these four
un-divided rows, and its real form is
\[
 C_0W^2+A_0Y^2+
                 (D_0+\tfrac12)(W+vX)^2+\tfrac12\Gamma^2\succ0.
 \tag{WN42}
\]
The coefficient-extraction substitution \(X=R/v,Z=\Gamma-sR/v\)
cancels all denominators; the asserted identity itself contains
no \(v^{-1}\), and includes literal \(v=0\).

For both reciprocal angles \(s,b=o(1)\), the full polynomial
belongs to the square ideal of \((W,X,Y,sZ,bZ)\), with real form
\[
 C_0W^2+A_0Y^2+D_0(X-sZ)^2+
                    \tfrac12X^2+\tfrac12(Y-bZ)^2\succ0.
 \tag{WN43}
\]
These finite identities include \(s=0\) and \(b=0\) exactly.

It remains to check that the local state charts were exhaustive.
Normalize a maximal state coordinate in the testing field.
At \(u,v=o(1)\), a unit \(X\) or \(Y\) gives positive full
residue immediately. Otherwise one of \(W,Z\) is a unit.
The exact sign gauges give \(u,v\geq0\).
If \(WZ\leq0\), the full-core sign identity SE9 supplies
separated nonnegative terms; all event packets are positive.
If \(WZ>0\), choose \(W=1,Z\) bounded, or \(Z=1,W=o(1)\).
These are exactly the states covered above; vanishing angles
or states are included by the same polynomial identities
(or by the nonpositive-sign case).
Positive pure-field homogeneous normalization and real/imaginary
state splitting now restore the original full parent:
\[
 F_c(\hbox{any independent projective test and state})
       \ \in\ \Sigma(K_c\otimes_{\mathbb R}K_t)^2 .
 \tag{WN44}
\]
This proves precisely the weak-\(n\) C2 and all-\(\beta\) C1
weak-\(\alpha\) source domains in the scope paragraph.
It supplies the missing local leaves usually called WB, WU
and CW. It does not, by itself, assert the global source
partition or the finite-lift conclusion.

\paragraph{Verification boundary.}
The focused checker reconstructs the entire parent, both graph
scalars, every first normal, the exact source balances, all
radical conductors, both exceptional normal completions and the
complete endpoint row identities. Source-order inequalities,
positive-unit residue arguments and the quantified cover are
proved above and in the expressly scoped WA source-entry lemma.
No missing historical proof, mixed inverse, or generic
SOS-cancellation assertion is used.

